\begin{table}[!htbp]
\centering\zihao{-4}\setlength{\tabcolsep}{4pt}
\begin{minipage}{\textwidth}\songti\normalsize
\subsection{公开方法适配与对比}\label{subsec:q2_public}
\QTwoParagraph 完成MMP与LTARP的内部比较及候选结构分析后，为考察LTARP与其他多模态建模路线的相对表现，本文进一步选取11种公开方法进行统一适配比较。各方法使用附件2的aligned-50输入与分类、回归双任务接口，三模态维度、数据划分、指标和完整输入选模标准保持一致，并按既有适配配置训练。表~\ref{tab:q2_public_overview}概括适配中的机制与规模。
\end{minipage}\par\vspace{10bp}
\caption{公开方法的统一接口适配概览}\label{tab:q2_public_overview}
\begin{tabularx}{.98\textwidth}{@{}>{\raggedright\arraybackslash}p{.17\textwidth}>{\raggedright\arraybackslash}p{.15\textwidth}>{\raggedright\arraybackslash}X r@{}}
\toprule
方法 & 类型 & 本文适配中的核心机制 & 参数量\\
\midrule
TFN\cite{q2tfn} & 张量融合 & 三路表示增广后的外积融合 & 4,840,670\\
LMF\cite{q2lmf} & 低秩融合 & 模态专属低秩因子融合 & 323,276\\
MFN\cite{q2mfn} & 时序记忆融合 & 三路递归编码、记忆注意力与门控记忆 & 261,764\\
MulT\cite{q2mult} & 跨模态注意力 & 六路定向跨模态注意力与模态内记忆 & 624,094\\
MISA\cite{q2misa} & 表征解耦 & 共享/特有表示、重构与差异约束 & 1,118,852\\
Self-MM\cite{q2selfmm} & 多任务辅助学习 & 融合预测与单模态伪目标辅助监督 & 82,119\\
MMIM\cite{q2mmim} & 互信息约束 & 模态间及融合--单模态的InfoNCE约束 & 163,204\\
\addlinespace
TFR-Net\cite{q2tfr} & 缺失重构 & 连续遮挡、时序上下文与潜在表示重构 & 415,620\\
MissModal\cite{q2missmodal} & 缺失表征对齐 & 几何对比、分布与情感语义对齐 & 163,204\\
M3S\cite{q2m3s} & 缺失元采样 & 低秩骨干、完整支持/缺失查询的一步一阶元更新 & 323,276\\
MMIN\cite{q2mmin} & 缺失模态推断 & 潜在表示推断、残差细化与循环约束 & 292,164\\
\bottomrule
\end{tabularx}
\par\vspace{10bp}\begin{minipage}{\textwidth}\songti\normalsize

\end{minipage}
\end{table}
